\clearpage
\section{Arbitrary contraction trees and deterministic store tracing}
\label{sec:copying-tracing}
The store rules of Figure~19 locate a store by following atom contraction
trees. This section proves two prerequisites for that matching operation.
First, every edge of an arbitrary finite acyclic contraction tree reaches
its unique root. Second, two successful atom traces in the same host agree
on the store they find. The second theorem does not assume that the host
is acyclic: stored references may form cycles.

Neither result is a definition of reference creation, reading, assignment,
or equality reduction. Nor do they prove that every atom in an arbitrary
host reaches a store. They establish structural termination for contraction
trees and uniqueness of successful store lookup, to be used by the future
machine rules. They also do not prove the copying rule's independence of
its chosen contraction-tree match.

\subsection{Exact traces and finite termination}
\termdef{exact-graph-trace}
For a relation $R$, an exact trace $a\mathrel{R^n}b$ consists of $n$ successive
steps. Its zero-step constructor ends where it starts; its step constructor
prepends one step. A terminal element has no outgoing $R$-step. Here $n$
counts graph adjacencies or atom hops, not machine transitions or the cost
relation used in contextual refinement.

Mapping every element of a trace along a relation-preserving function
preserves its exact length, by induction on the trace constructors. A
bijection which preserves and reflects steps also reflects traces: map
back through its inverse, use step reflection, and apply the two inverse
equations to the endpoints. These lemmas will transfer store lookup across
graph renaming and depth padding.

\lean{Graph/FiniteTrace}{7}{9}{FiniteTrace.Trace}

\lean{Graph/FiniteTrace}{11}{14}{FiniteTrace.Terminal}

\lean{Graph/FiniteTrace}{17}{22}{FiniteTrace.map}

\lean{Graph/FiniteTrace}{25}{37}{FiniteTrace.map_iff}

If $R$ is functional, two traces starting at the same element and ending
at terminal elements have equal endpoints and equal lengths. Induct on the
first trace. A zero-step trace cannot compete with a positive trace, because
its endpoint is terminal. If both traces take a step, functionality identifies
the next elements; induction identifies the remaining endpoints and lengths.

For an arbitrary relation, bounded exploration either reaches a terminal
or produces the requested number of steps. At bound zero take the reflexive
trace. At a positive bound, stop if the current element is terminal;
otherwise classical reasoning supplies a successor. Apply the induction
hypothesis there and prepend the chosen step. This is a proof of existence,
not an executable decision procedure for arbitrary relations.

\lean{Graph/FiniteTrace}{40}{55}{FiniteTrace.terminal_unique}

\lean{Graph/FiniteTrace}{58}{75}{FiniteTrace.terminal_or_long}

An adjacency trace of length $n$ in a graph gives a path of $n+1$ edges:
induction prepends the starting edge at each step. If every path has no
repeated edge, its length is at most the length of the graph's finite edge
enumeration, since every path edge occurs in that enumeration. Repeated
entries in the enumeration are harmless. Consequently the trace length is
strictly less than the enumeration length.

Apply bounded exploration with the enumeration length as its bound. The
alternative producing that many adjacency steps contradicts this strict
inequality. Thus every edge in a finite edge-acyclic graph reaches a terminal
edge. This proof does not assume a rank, a balanced tree, or a chosen
canonical presentation.

\lean{Graph/FiniteTrace}{78}{84}{FiniteTrace.trace_path}

\lean{Graph/FiniteTrace}{87}{94}{FiniteTrace.trace_bound}

\lean{Graph/FiniteTrace}{97}{102}{FiniteTrace.exists_terminal_trace}

If an edge has exactly one target vertex, two adjacent successors must
share that same vertex in their source lists. Source-incidence uniqueness
then makes the two successors equal. This local fact supplies functionality
both for contraction trees and for atom tracing in larger hosts.

\lean{Graph/FiniteTrace}{105}{117}{FiniteTrace.adjacent_functional_at}

\subsection{Arbitrary contraction trees on concrete graphs}
\termdef{arbitrary-contraction-tree}
For sharing sort $s$ and $m$ exposed leaves, a contraction tree is a concrete
finite graph $s^m\to s$ whose edges are all $s$-contractions or $s$-weakenings
and whose paths have no repeated edge. The structure does not assume
connectedness to the output; that is proved below. It does not require a
weakening edge: the single binary contraction is a valid tree. This is
broader than the canonical trees, which have exactly one weakening.

The structure records the finite graph, its label restriction, and the
acyclicity proof. Both permitted labels have a single target by their
signatures, so every tree edge has exactly one target vertex. The generic
single-target lemma therefore proves that forward adjacency is functional.

\lean{Copying/ContractionTree}{11}{17}{ContractionTree.ContractionTree}

\lean{Copying/ContractionTree}{20}{22}{ContractionTree.target_single}

\lean{Copying/ContractionTree}{24}{26}{ContractionTree.adjacent_functional}

An edge is terminal exactly when its target is the sole external output.
For the forward implication, inspect its unique target vertex. If that
vertex were not an output, source coverage would supply a successor edge,
contradicting terminality. For the reverse implication, the unique target
is an output, and outputs never occur in source lists, so no successor
exists.

Target coverage supplies an edge producing the external output. Define the
root to be that edge, using classical choice. Its output incidence follows
from the choice specification, and the preceding equivalence makes it
terminal. Any other terminal edge produces the same output, so target
uniqueness identifies it with the root. Thus the chosen name has no effect
on root uniqueness.

\lean{Copying/ContractionTree}{29}{60}{ContractionTree.terminal_iff_output}

\lean{Copying/ContractionTree}{63}{64}{ContractionTree.root}

\lean{Copying/ContractionTree}{66}{67}{ContractionTree.root_output}

\lean{Copying/ContractionTree}{69}{70}{ContractionTree.root_terminal}

\lean{Copying/ContractionTree}{72}{73}{ContractionTree.terminal_eq_root}

Finite acyclic termination supplies a terminal trace from every edge.
Its endpoint must be the root. Functionality additionally shows that any
two terminal traces from that edge have the same endpoint and length.
For an exposed input, source coverage first supplies its consuming edge;
then use that edge's root trace. The same argument applies to every vertex
other than the external output. This derives the required connection of
the entire tree to its root, rather than placing connectivity in the
structure's assumptions.

\lean{Copying/ContractionTree}{76}{80}{ContractionTree.reaches_root}

\lean{Copying/ContractionTree}{83}{86}{ContractionTree.root_trace_unique}

\lean{Copying/ContractionTree}{89}{93}{ContractionTree.input_reaches_root}

\lean{Copying/ContractionTree}{96}{100}{ContractionTree.vertex_reaches_root}

A graph isomorphism transports the tree property. Labels are preserved,
and a path in the renamed graph maps back to a path in the original tree.
The latter has no repeated edge; equality of positions in the former would
map to equality of positions in the latter, so the former is also without
repetitions. Finally, each canonical tree is an instance: its earlier label
classification supplies the sharing-label restriction, and its rank proof
supplies acyclicity.

\lean{Copying/ContractionTree}{103}{115}{ContractionTree.transport}

\lean{Copying/ContractionTree}{118}{124}{ContractionTree.canonical}

\subsection{Atom tracing in arbitrary hosts}
\termdef{atom-store-resolution}
An atom hop may leave an instance edge, an atom contraction, or an atom
weakening and follows adjacency to the next edge. These three labels have
one atom output. A successful store resolution is an exact sequence of
such hops whose endpoint is labelled as a store. Unlike unrestricted graph
adjacency, this relation stops at a store and does not follow its term output.

The forward-label predicate lists the three allowed labels at any box
allowance. A hop is that predicate together with adjacency. The output
signature of every forward label is the singleton atom interface, by cases
on the predicate. Hence the target incidence list has length one, and
single-target adjacency functionality proves uniqueness of the next hop.
No hypothesis about cycles elsewhere in the host is used.

\lean{Copying/StoreTrace}{13}{16}{StoreTrace.ForwardLabel}

\lean{Copying/StoreTrace}{19}{20}{StoreTrace.Hop}

\lean{Copying/StoreTrace}{22}{24}{StoreTrace.forward_output}

\lean{Copying/StoreTrace}{27}{31}{StoreTrace.hop_functional}

The common vertex of a hop lies in the forward edge's target list. Target
typing makes its label an atom; source typing therefore puts an atom in the
successor edge's input signature. In a hole-free host this forces the next
edge to be an atom contraction or a store. Inspect the alphabet: operations,
instances, and focuses have term inputs; boxes have thunk inputs; weakenings
have no inputs; term contractions have term inputs; holes are excluded by
the global hole-free predicate. Thus the permitted successors are proved
from typing and the alphabet, rather than assumed as a matching convention.

\lean{Copying/StoreTrace}{34}{43}{StoreTrace.hop_atom_input}

\lean{Copying/StoreTrace}{46}{63}{StoreTrace.hop_successor_kind}

A store is terminal for atom hopping because its output is a term, whereas
all forward labels output an atom. A successful resolution records the
trace and the final store label. Apply the functional-relation terminal
uniqueness theorem to two successful resolutions: both store endpoints and
hop counts agree. Apply this result separately to two starting atoms to
show that testing whether their resolved store edges are equal is independent
of the chosen successful traces. This establishes the lookup premise needed
by atom equality; it does not yet prove the equality compute rule.

\lean{Copying/StoreTrace}{66}{71}{StoreTrace.store_terminal}

\lean{Copying/StoreTrace}{74}{75}{StoreTrace.Resolves}

\lean{Copying/StoreTrace}{78}{81}{StoreTrace.resolves_unique}

\lean{Copying/StoreTrace}{84}{90}{StoreTrace.same_store_iff}

A graph isomorphism preserves labels and preserves and reflects adjacency,
so it preserves and reflects atom hops. The exact-trace bijection lemma and
preservation of the final store label then give equivalence of successful
resolutions under renaming, including equality of the hop counts.

\lean{Copying/StoreTrace}{93}{95}{StoreTrace.hop_iso_iff}

\lean{Copying/StoreTrace}{98}{101}{StoreTrace.resolves_iso_iff}

Depth padding leaves the edge carrier unchanged; an explicit bijection
exposes this fact despite the depth-indexed type. At depth zero, lifted
primitive labels become primitive labels at depth one. At positive depth,
primitive labels remain primitives and box labels remain boxes. Case
analysis therefore preserves and reflects the forward-label predicate and
the store-label predicate. Incidence lists are unchanged. These facts give
hop equivalence, and exact-trace transport gives resolution equivalence
at the larger allowance. Padding creates no new store lookup behaviour.

\lean{Copying/StoreTrace}{104}{107}{StoreTrace.paddingEdges}

\lean{Copying/StoreTrace}{110}{121}{StoreTrace.hop_raise_iff}

\lean{Copying/StoreTrace}{124}{135}{StoreTrace.resolves_raise_iff}

\subsection{Regressions: no weakening and a cyclic store}

The weakening-free binary tree is the corolla consisting of one contraction.
Its label condition is immediate. It has no adjacency because its targets
are output ports and its sources are input ports. A constant rank therefore
vacuously increases on adjacency and proves path acyclicity. The generic
tree theorem applies to every edge of this tree and proves root reachability.
This checks that the arbitrary-tree definition has not silently inherited
the canonical tree's extra weakening.

\lean{Copying/TracingExamples}{9}{17}{TracingExamples.binaryTree}

\lean{Copying/TracingExamples}{20}{22}{TracingExamples.binaryTree_reaches}

The cyclic-store graph has two internal vertices, one term and one atom,
and two edges. An instance runs from the term vertex to the atom vertex;
a store runs from the atom vertex back to the term vertex. It has empty
external interfaces. Boolean enumeration supplies finiteness, and direct
case analysis verifies all typing, incidence uniqueness, and coverage
fields. The path consisting of instance, store, instance witnesses a genuine
cycle, so global acyclicity is unavailable in this example.

\lean{Copying/TracingExamples}{25}{58}{TracingExamples.cyclicStore}

\lean{Copying/TracingExamples}{61}{64}{TracingExamples.cyclicStore_cycle}

Nevertheless, its instance resolves to its store in exactly one atom hop:
the shared atom vertex supplies adjacency, and the endpoint has the store
label. Any successful resolution of that instance must have the same store
and length by the general uniqueness theorem. Both labels are different
from every hole label, giving global hole freedom at every depth allowance.
This checks that the tracing relation stops before following a stored value
back into the cycle. It is a graph-level regression, not a machine execution
or an allocation/read/write test.

\lean{Copying/TracingExamples}{67}{69}{TracingExamples.cyclicStore_resolves}

\lean{Copying/TracingExamples}{72}{74}{TracingExamples.cyclicStore_unique}

\lean{Copying/TracingExamples}{77}{79}{TracingExamples.cyclicStore_holeFree}

All new declarations pass the full build and axiom audit with only Lean's
standard axioms. Arbitrary contraction-tree construction and root tracing
are completed prerequisites. Copyable and stable hypernets are now defined
in Section~\ref{sec:copying-values}. Context translation and the translation laws,
actual copying and store transitions, their full matching and
locality proofs, the templates, and Proposition~12.2 remain open.
